\chapter*{Abstract}
\addcontentsline{toc}{chapter}{Abstract}
Automotive zonal electrical/electronic architectures place safety-critical control
traffic, sensor streams and infotainment on a single Ethernet backbone, and rely on
Time-Sensitive Networking (TSN) to keep that traffic deterministic. Determinism
depends on the generalised Precision Time Protocol (gPTP, IEEE 802.1AS), which was
standardised without security controls. An attacker who manipulates synchronisation
can therefore degrade timing guarantees without altering packet content or traffic
volume, which places the attack outside the detection surface of conventional
content- and volume-based intrusion detection.

This project develops a federated intrusion detection system that operates on
TSN protocol-native signals rather than generic traffic statistics. Zone controllers
act as federated clients, each extracting timing features from locally observed gPTP
state and training a shared model architecture; a central vehicle compute unit
aggregates model parameters without receiving raw feature records.

At the mid-report stage, a four-zone automotive TSN topology has been implemented in
OMNeT++/INET with gPTP synchronisation, and three attack configurations have been
instrumented. Fifteen seeded simulation runs (three benign and twelve delay-attack runs covering
four target zones at three seeds each) have produced 225 recorded vectors,
converted by a Python extraction pipeline into 4{,}740 labelled feature windows.
Delay-attack injection has been validated against the expected theoretical signature.
A cross-zone residual feature has been added and a centralized detector trained,
establishing the project's central premise ahead of schedule: features computed within
a single zone detect the attack no better than a majority-class predictor
(macro-F1 0.48--0.54 against a floor of 0.47), whereas two cross-zone residual features
reach 0.95--0.98 macro-F1 with detection two synchronisation intervals after onset.
Two negative results are reported: oscillator drift manipulation is not realisable as
a runtime attack under the simulation framework, and no fixed threshold, neither a
maximum nor a percentile, separates attack from benign behaviour, because the
synchronisation convergence transient produces residuals 2.7 times larger than the
attack does. Remaining work for FYP-2 comprises zone heterogeneity, the federated
training pipeline, and the comparative evaluation.
